%!TEX program = xelatex
\title          {Logic in Computer Science - Assignment 3}
\author         {李鹏达 10225101460}
\date{}

\documentclass{article}
\usepackage{amsmath}
\usepackage{../boxproof}
\usepackage{a4wide}
\usepackage{ctex}
% \usepackage{daymonthyear}

\def\meta#1{\mbox{$\langle\hbox{#1}\rangle$}}
\def\macrowitharg#1#2{{\tt\string#1\bra\meta{#2}\ket}}

{\escapechar-1 \xdef\bra{\string\{}\xdef\ket{\string\}}}

\showboxbreadth 999
\showboxdepth 999
\tracingoutput 1


\let\imp\to

\def\premise{\mathrm{premise}}
\def\assumption{\mathrm{assumption}}
\def\MT{\mathrm{MT\ }}
\def\LEM{\mathrm{LEM}}
\def\intro{\mathrm{i\ }}
\def\elim{\mathrm{e\ }}
\def\introa{\mathrm{i_1\ }}
\def\elima{\mathrm{e_1\ }}
\def\introb{\mathrm{i_2\ }}
\def\elimb{\mathrm{e_2\ }}

\def\lt{<}
\def\eqdef{=}

\def\eps{\mathrel{\epsilon}}
\def\biimplies{\leftrightarrow}
\def\flt#1{\mathrel{{#1}^\flat}}
\def\setof#1{{\left\{{#1}\right\}}}
\let\implies\to
\def\KK{{\mathsf K}}
\let\squashmuskip\relax


\def\pre{\premise}
\def\ass{\assumption}
\def\andea{\land\elima}
\def\andeb{\land\elimb}
\def\andi{\land\intro}
\def\toe{\to\elim}
\def\toi{\to\intro}
\def\ore{\lor\elim}
\def\ore{\lor\elim}
\def\oria{\lor\introa}
\def\orib{\lor\introb}
\def\nege{\neg\elim}
\def\negi{\neg\intro}
\def\bote{\bot\elim}

\def\tt{\!-\!}

%=======================================================================
\begin{document}
\maketitle
%=======================================================================
\noindent 1. For completeness proof, complete the proof of the following cases:
\\[1em]
(1) $\Phi_1 \land \Phi_2 \vdash \Phi_1 \lor \Phi_2$
$$\begin{proofbox}
\: \Phi_1 \land \Phi_2 \= \pre \\
\: \Phi_1 \= \andea 1 \\
\: \Phi_1 \lor \Phi_2 \= \oria 2 \\
\end{proofbox}$$
\\[1em]
(2) $\Phi_1 \land \neg \Phi_2 \vdash \Phi_1 \lor \Phi_2$
$$\begin{proofbox}
\: \Phi_1 \land \neg\Phi_2 \= \pre \\
\: \Phi_1 \= \andea 1 \\
\: \Phi_1 \lor \Phi_2 \= \oria 2 \\
\end{proofbox}$$
\\[1em]
(3) $\neg \Phi_1 \land \Phi_2 \vdash \Phi_1 \lor \Phi_2$
$$\begin{proofbox}
\: \neg\Phi_1 \land \Phi_2 \= \pre \\
\: \Phi_2 \= \andeb 1 \\
\: \Phi_1 \lor \Phi_2 \= \orib 2 \\
\end{proofbox}$$
\\[1em]
(4) $\neg \Phi_1 \land \neg \Phi_2 \vdash \neg(\Phi_1 \lor \Phi_2)$
$$\begin{proofbox}
\: \neg\Phi_1 \land \neg\Phi_2 \= \pre \\
\: \neg \Phi_1 \= \andea 1 \\
\: \neg \Phi_2 \= \andeb 1 \\
\[
    \: \Phi_1 \lor \Phi_2 \= \ass \\
    \[
        \: \Phi_1 \= \oria 4 \\
        \: \bot \= \nege 2, 5 \\
    \]
    \[
        \: \Phi_2 \= \orib 4 \\
        \: \bot \= \nege 3, 7 \\
    \]
    \: \bot \= \ore 4,5\tt6,7\tt8 \\
\]
\: \neg(\Phi_1 \lor \Phi_2) \= \negi 4\tt9 \\
\end{proofbox}$$
\\[1em]
(5) $\Phi_1 \land \Phi_2 \vdash\Phi_1 \to \Phi_2$
$$\begin{proofbox}
\: \Phi_1 \land \Phi_2 \= \pre \\
\[
    \: \Phi_1 \= \ass \\
    \: \Phi_2 \= \andeb 1 \\
\]
\: \Phi_1 \to \Phi_2 \= \toi 2\tt3 \\
\end{proofbox}$$
\\[1em]
(6) $\Phi_1 \land \neg \Phi_2 \vdash \neg(\Phi_1 \to \Phi_2)$
$$\begin{proofbox}
\: \Phi_1 \land \neg\Phi_2 \= \pre \\
\: \Phi_1 \= \andea 1 \\
\: \neg \Phi_2 \= \andeb 1 \\
\[
    \: \Phi_1 \to \Phi_2 \= \ass \\
    \: \Phi_2 \= \toe 4, 2 \\
    \: \bot \= \nege 3, 5 \\
\]
\: \neg(\Phi_1 \to \Phi_2) \= \negi 4\tt6 \\
\end{proofbox}$$
\\[1em]
(7)  $\neg \Phi_1 \land \Phi_2 \vdash\Phi_1 \to \Phi_2$
$$\begin{proofbox}
\: \neg \Phi_1 \land \Phi_2 \= \pre \\
\[
    \: \Phi_1 \= \ass \\
    \: \Phi_2 \= \andeb 1 \\
\]
\: \Phi_1 \to \Phi_2 \= \toi 2\tt3 \\
\end{proofbox}$$
\\[1em]
(8) $\neg \Phi_1 \land \neg \Phi_2 \vdash \Phi_1 \to \Phi_2$
$$\begin{proofbox}
\: \neg \Phi_1 \land \neg \Phi_2 \= \pre \\
\[
    \: \Phi_1 \= \ass \\
    \: \neg \Phi_1 \= \andea 1 \\
    \: \bot \= \nege 2, 3 \\
    \: \Phi_2 \= \bote 4 \\
\]
\: \Phi_1 \to \Phi_2 \= \toi 2\tt5 \\
\end{proofbox}$$
\newpage
\noindent 2. Convert the following formula into CNF
$$
\neg\left(r \land \left(\neg \left(q \to \left(\neg p \to \left(q \land r\right)\right)\right)\right)\right)
$$

\noindent\textbf{Solution:}
\\[2em]
Let $\Phi = \neg\left(r \land \left(\neg \left(q \to \left(\neg p \to \left(q \land r\right)\right)\right)\right)\right)$.
\newcommand{\cnf}[1]{\texttt{CNF($#1$)}}
\newcommand{\nnf}[1]{\texttt{NNF($#1$)}}
\newcommand{\impf}[1]{\texttt{IMPL\_FREE($#1$)}}
\newcommand{\distr}[2]{\texttt{DISTR($#1,#2$)}}
\begin{align*}
& \impf{\Phi}\\
&= \impf{\neg\left(r \land \left(\neg \left(q \to \left(\neg p \to \left(q \land r\right)\right)\right)\right)\right)}\\
&= \neg(\impf{r \land \left(\neg \left(q \to \left(\neg p \to \left(q \land r\right)\right)\right)\right)})\\
&= \neg(\impf{r} \land \impf{\neg \left(q \to \left(\neg p \to \left(q \land r\right)\right)\right)})\\
&= \neg(r \land \neg(\impf{q \to \left(\neg p \to \left(q \land r\right)\right)}))\\
&= \neg(r \land \neg(\neg \impf{q} \lor \impf{\neg p \to \left(q \land r\right)}))\\
&= \neg(r \land \neg(\neg q \lor (\neg \impf{p} \lor \impf{q \land r})))\\
&= \neg(r \land \neg(\neg q \lor (\neg p \lor (\impf{q} \land \impf{r}))))\\
&= \neg(r \land \neg(\neg q \lor (\neg p \lor (q \land r))))\\
\end{align*}
Let $\Phi_1 = \neg(r \land \neg(\neg q \lor (\neg p \lor (q \land r))))$.
\begin{align*}
& \nnf{\Phi_1} \\
&= \nnf{\neg(r \land \neg(\neg q \lor (\neg p \lor (q \land r))))} \\
&= \nnf{\neg r \lor \neg\neg(\neg q \lor (\neg p \lor (q \land r)))} \\
&= \nnf{\neg r} \lor \nnf{\neg\neg(\neg q \lor (\neg p \lor (q \land r)))} \\
&= \neg r \lor \nnf{\neg q \lor (\neg p \lor (q \land r))} \\
&= \neg r \lor \nnf{\neg q} \lor \nnf{\neg p \lor (q \land r)} \\
&= \neg r \lor \neg q \lor \nnf{\neg p} \lor \nnf{q \land r} \\
&= \neg r \lor \neg q \lor \nnf{\neg p} \lor (\nnf{q} \land \nnf{r}) \\
&= \neg r \lor \neg q \lor \neg p \lor (q \land r) \\
\end{align*}
Let $\Phi_2 = \neg r \lor \neg q \lor \neg p \lor (q \land r)$.
\begin{align*}
& \cnf{\Phi_2} \\
&= \cnf{\neg r \lor \neg q \lor \neg p \lor (q \land r)} \\
&= \distr{\cnf{\neg r \lor \neg q}}{\cnf{\neg p \lor (q \land r)}} \\
&= \distr{\distr{\cnf{\neg r}}{\cnf{\neg q}}}{\distr{\cnf{\neg p}}{\cnf{q \land r}}} \\
&= \distr{\distr{{\neg r}}{{\neg q}}}{\distr{{\neg p}}{{\cnf{q} \land \cnf{r}}}} \\
&= \distr{\distr{{\neg r}}{{\neg q}}}{\distr{{\neg p}}{{q \land r}}} \\
&= \distr{\neg r \lor \neg q}{\distr{\neg p}{q}\land \distr{\neg p}{r}} \\
&= \distr{\neg r \lor \neg q}{({\neg p}\lor{q})\land ({\neg p}\lor{r})} \\
&= \distr{\neg r \lor \neg q}{{\neg p}\lor{q}} \land \distr{\neg r \lor \neg q}{{\neg p}\lor{r}} \\
&= (\neg r \lor \neg q \lor \neg p \lor q) \land (\neg r \lor \neg q \lor \neg p \lor r)
\end{align*}
Thus, the CNF of the formula is
$$(\neg r \lor \neg q \lor \neg p \lor q) \land (\neg r \lor \neg q \lor \neg p \lor r)$$
\end{document}

